\section{Processor}
The processor comprises two major components: the \textbf{datapath} and the \textbf{control}.
\subsection{Datapath}
The datapath is the \textit{collection of components that process data}. It
performs the arithmetic, logical and memory operations.
\begin{remark}
  This processor design supports a subset of the core MIPS ISA:
  \begin{itemize}
    \item Arithmetic and Logical operations:
      \verb|add, sub, and, or, addi, slt, sll, srl|
    \item Data Transfer Instructions
      \verb|lw, sw|
    \item Branches
      \verb|beq, bne|
    \item Jumps
      \verb|j|
  \end{itemize}
  \verb|andi, ori| is not supported because sign extension is done on immediate
  values.
\end{remark}
\subsubsection{Instruction Execution Cycle}
The Instruction Execution Cycle comprises the following steps:
\begin{figure}[h!]
  \centering
  \includegraphics[width=0.5\textwidth]{IEC.jpeg}
\end{figure}
\paragraph{Instruction Fetch Stage}: Uses PC to fetch the instruction from
\textbf{instruction memory}, and then is incremented by $4$ to get the address of the next
instruction. Output instruction is then fed to the next stage.

\begin{definition}[Instruction Memory]
  The instruction memory is a sequential circuit that as an internal state
  storing the instructions. It also comprises a clock signal.
\end{definition}

\paragraph{Clock Period}
The clock signal ensures that we do not read and write the PC at the same time,
by ensuring that the PC is read during the first half of the clock period, and
updated at the \textbf{next rising clock edge}.\\
In general, the clock period allows us to read, compute, and write all within a
single clock period, and allow us to synchronise between each phase of the
instruction execution.


\paragraph{Instruction Decode/Operand Fetch Stage}: Gathers the data from the
instruction fields. It determines the instruction type and field lengths, and
then read from all necessary registers. The output operation and operands is
then fed to the next stage.

\begin{definition}[Register File]
  The register file is a collection of 32 registers that can be read/written by
  specifying the register number. Per instruction, a maximum of two registers
  can be read and one register can be written.
\end{definition}

\paragraph{Instruction ALU Stage}: Executes an operation specified by
ALUControl with the two input values. It outputs \verb|ALUresult| and
\verb|is0?|. The \verb|is0?| bit is used to generate the PCSrc Control signal
that decides whether there is a jump.

\begin{figure}[h!]
  \centering
  \includegraphics[width=0.25\textwidth]{aluControl.jpeg}
\end{figure}

\paragraph{Memory Stage}: For load and store instructions, we use the memory
address calculated by the ALU stage to read or write data memory. The output value
is passed on to the RegWrite stage.

\begin{definition}[Data Memory]
  The data memory is a storage element for the data of a program., taking in a
  memory address, data to be written and outputs the data read from memory for
  load instructions.
\end{definition}

\paragraph{RegWrite Stage}: For most instructions, we write the result of computation
into a register (except for stores, branches and jumps).

\subsection{Control}
The control prescribes the processes of the datapath, memory and I/O devices
according to program instructions, through control signals. From the datapath,
we are required to generate the control signals:

\begin{figure}[h!]
  \centering
  \includegraphics[width=0.7\textwidth]{controlSignals.jpeg}
\end{figure}

\paragraph{ALUControl Design}
Considering the 1-bit ALU, the ALUControl signal indicates whether to invert
each of the input, and specify the final operation performed on the inputs, as
shown in the figure.

\begin{figure}[h!]
  \centering
  \includegraphics[width=0.7\textwidth]{aluControl2.jpeg}
\end{figure}

We can design the ALUControl signal which depends on the opcode and function
code fields with a \textit{multilevel decoding approach}, by generating an
intermediate signal \verb|ALUop| from the 6-bit opcode, and then obtain
\verb|ALUcontrol| from \verb|ALUop| and funct.

\begin{remark}
  The multilevel decoding approach simplifies the design process, reduces the
  size of the main controller, and potentially speeds up the circuit.
\end{remark}

\begin{figure}[h!]
  \centering
  \includegraphics[width=0.5\textwidth]{aluControlCircuit.png}
\end{figure}

\begin{figure*}[h!]
  \centering
  \includegraphics[width=0.65\textwidth]{aluControlOverview.jpeg}
  \captionsetup{labelformat=empty}
  \caption{ALUop $\to$ ALUcontrol}
\end{figure*}

\paragraph{Control Design}
With the required outputs of the control, we now design the control circuit using the
opcode.
\begin{figure}[h!]
  \centering
  \includegraphics[width=0.4\textwidth]{controlCircuit.png}
\end{figure}
\begin{figure}[h!]
  \centering
  \includegraphics[width=0.7\textwidth]{control.jpeg}
  \includegraphics[width=0.4\textwidth]{control2.jpeg}
  \captionsetup{labelformat=empty}
  \caption{opcode $\to$ Control Signals}
\end{figure}

\subsection{Instruction Execution}
Each instruction consists of three phases: \textit{Read} contents of one or
more storage elements, Perform \textit{computations} through some combinational
logic, \textit{Write} results to one or more storage elements.

Each instruction is performed within a single \textbf{clock period}. The clock
period is used to prevent concurrent reading and writing to the storage
elements.
\begin{figure}[h!]
  \centering
  \includegraphics[width=0.5\textwidth]{clockPeriod.png}
\end{figure}

The clock period is computed based on the instructions to be executed, and has
multiple possible implementations.

\subsubsection{Single Cycle Implementation}
In Single Cycle Implementation, we calculate the clock cycle time based on the
slowest instruction between \textit{ALU}, \textit{lw}, \textit{sw},
\textit{beq} instructions.
\begin{example}
  Assuming neglligible delays, we have the total time for each instruction:
  \begin{figure}[h!]
    \centering
    \includegraphics[width=0.6\textwidth]{singleCycle.png}
  \end{figure}
\end{example}
Implementing a single cycle for all instructions, each instruction will take as
much time as the slowest one.
  
\subsubsection{Multicycle Implementation}
In multicycle implementation, we break up instructions into \textit{execution
steps}, each taking a single clock cycle.
\begin{itemize}
  \item Instruction Fetch
  \item Instruction Decode and Register Read
  \item ALU Operation
  \item Memory Read/Write
  \item Register Write
\end{itemize}
The cycle time is hence much shorter, and each instruction takes a variable
number of clock cycles.

\subsubsection{Pipelining}
In addition to breaking up instructions into execution steps, we may
\textit{allow different instructions to be in different execution steps
simultaneously}. This is known as \textbf{pipelining}.

\newpage
\begin{sideways}[h!]
  \centering
  \includegraphics[width=1.35\textwidth]{processor.jpeg}
\end{sideways}
\newpage
